% !PRINTED 547
\weberpage{547}
%<<CONT
and has no others. In this lies the proof of the following theorem:

3. \gesperrt{A number $\omega$ in $\mathfrak{o}$ is then and only then congruent to a rational number with respect to the modulus $\pi$, when the condition}
\[ \omega^p \equiv \omega \pmod{\pi} \]
\gesperrt{is satisfied.}

If one observes further that the polynomial coefficients in the $p$th power of a polynomial are all divisible by $p$, with the exception of those which belong to the $p$th powers of the individual terms of the polynomial, there results the following further theorem, which relates to integral functions in $\Omega$ of arbitrary variables $x, y, \ldots$:

4. \gesperrt{If $\psi(x,y,\ldots)$ is an integral function of the variables $x, y, \ldots$ with integral coefficients from $\Omega$, then the holding of the congruence}
\[ [\psi(x,y,\ldots)]^p \equiv \psi(x^p, y^p, \ldots) \pmod{\pi} \]
\gesperrt{is the necessary and sufficient condition that all coefficients of $\psi$ are congruent to integral rational numbers with respect to the modulus $\pi$.}

\begin{center}
\S.\,151.\\
\gesperrt{The Dedekind Ideals.}
\end{center}

\gesperrt{Dedekind} founds the theory of algebraic numbers, in the papers already mentioned above, on the concept of the \gesperrt{ideal}.

We now wish to prove that the theory of ideals agrees in essence with the theory of functionals, by showing how the transition from the one to the other can be effected.

The system of all integers of an algebraic number field $\Omega$ shall, as above, be denoted by $\mathfrak{o}$\footnote{Cf. \gesperrt{Dirichlet-Dedekind}, Vorlesungen über Zahlentheorie, in \S.\,167 of the third, \S.\,177 of the fourth edition.}. A system of numbers $\mathfrak{a}$ contained in $\mathfrak{o}$ is called an \gesperrt{ideal} if it satisfies the two requirements:
%CONT>>
